\begin{abstract}
Rad daje retrospektivnu analizu razvoja simetrične kriptografije od
Data Encryption Standarda (DES), preko Triple DES-a (TDEA), do Advanced
Encryption Standarda (AES), sa posebnim osvrtom na telekomunikacione i
mrežne sisteme. Analiziraju se dokumentovana očekivanja iz perioda izbora
i početne standardizacije AES-a i porede sa razvojem standarda,
implementacija i protokola do 2026. godine. Eksperimentalni deo uspostavlja
validiran i ponovljiv metodološki okvir za mikrobenchmarking
kriptografskih implementacija, sa jasnim razdvajanjem softverske,
hardverski akcelerovane, GPU resident i end-to-end putanje. Rezultati
pokazuju da AES-NI značajno smanjuje praktičnu cenu AES-a na savremenom
CPU-u, dok CUDA AES-128 ostvaruje višu izolovanu resident propusnost za
veće ulaze, ali u testiranom sinhronom host--device pipeline-u ne
nadmašuje AES-NI referencu. AES-GCM merenja dodatno pokazuju značaj
posmatranja kompletne autentifikovane obrade umesto same blokovske
transformacije. Demonstracije redukovanog prostora ključeva i avalanche
efekta koriste se kao ilustrativni, a ne sigurnosno dokazni eksperimenti.
Retrospektiva pokazuje da je AES izdržao test vremena kao dugovečna,
fleksibilna i efikasno implementabilna primitiva, dok su se praktični
izazovi pomerili ka implementacionoj bezbednosti, autentifikovanoj
enkripciji, upravljanju nonce vrednostima, ograničenim uređajima i
heterogenim računarskim arhitekturama.

\noindent\textbf{Ključne reči:} simetrična kriptografija; DES; TDEA; AES;
AES-NI; GPU akceleracija; autentifikovana enkripcija.
\end{abstract}


\begin{otherlanguage}{english}
\begin{abstract}
This paper presents a retrospective analysis of the evolution of symmetric
cryptography from the Data Encryption Standard (DES), through Triple DES
(TDEA), to the Advanced Encryption Standard (AES), with particular emphasis
on telecommunications and networked systems. Documented expectations from
the AES selection and early standardization period are compared with the
subsequent evolution of standards, implementations, and protocols through
2026. The experimental part establishes a validated and reproducible
methodological framework for microbenchmarking cryptographic
implementations, with explicit separation of software, hardware-accelerated,
GPU-resident, and end-to-end execution paths. The results show that AES-NI
substantially reduces the practical cost of AES on a modern CPU, while
CUDA AES-128 achieves higher isolated resident throughput for larger
inputs but does not outperform the AES-NI reference in the tested
synchronous host--device pipeline. AES-GCM measurements further demonstrate
the importance of evaluating complete authenticated processing rather than
block-cipher throughput alone. Reduced-keyspace search and avalanche-effect
experiments are used as illustrative measurements rather than proofs of
security. The retrospective indicates that AES has withstood the test of
time as a long-lived, flexible, and efficiently implementable primitive,
while practical challenges have increasingly shifted toward implementation
security, authenticated encryption, nonce management, constrained devices,
and heterogeneous computing architectures.

\noindent\textbf{Keywords:} symmetric cryptography; DES; TDEA; AES;
AES-NI; GPU acceleration; authenticated encryption.
\end{abstract}
\end{otherlanguage}